\section{Self-attention from scratch}
\label{ch:self-attention}

\newthought{If the transformer block} in Section~\ref{ch:transformer} has a
single piece worth understanding cell-by-cell, it is the self-attention
sub-layer. This section writes it out, with all the matrix shapes.

\subsection{The setup}

Take the input to one block: a tensor \(X \in \R^{n \times d}\), where \(n\)
is the number of tokens and \(d = \dmodel\) is the hidden dimension. Each row
\(x_t\) is the current representation of token \(t\).

We are going to produce an output of the same shape: a new representation
for each token that has incorporated information from earlier tokens.

\subsection{Queries, keys, and values}

Project the input into three different spaces:
\begin{equation}
  Q \;=\; X W_Q, \quad
  K \;=\; X W_K, \quad
  V \;=\; X W_V,
\end{equation}
with three learned matrices \(W_Q, W_K, W_V \in \R^{d \times \dhead}\). Each
of \(Q, K, V\) is \([n, \dhead]\).\sidenote{The intuition. \emph{Query} =
``what am I looking for?''. \emph{Key} = ``what do I advertise as being
about?''. \emph{Value} = ``what do I actually carry that you'd want to
read?''. The attention weight from token \(t\) to token \(s\) is high when
\(t\)'s query and \(s\)'s key align; the value \(v_s\) is what gets copied
into \(t\)'s next representation.}

\subsection{The attention matrix}

Form pairwise scores via dot product, scale by \(\sqrt{\dhead}\) so the
softmax does not saturate, and softmax row-wise:
\begin{equation}
  A \;=\; \softmax\!\left( \frac{Q K\T}{\sqrt{\dhead}} \right)
  \;\in\; \R^{n \times n}.
  \label{eq:attn-mat}
\end{equation}
\(A_{t,s}\) is ``how much should token \(t\) attend to token \(s\)''. The
softmax is taken over rows, so each row sums to one. Every token spreads
its attention budget over all other tokens.\sidenote{The scaling
\(1 / \sqrt{\dhead}\) is not a cosmetic choice. Without it, dot products
grow with \(\sqrt{\dhead}\) in expectation, the softmax saturates, and
gradients vanish. \citet{vaswani2017attention} document this in their
section 3.2.1.}

\subsection{The output}

Use \(A\) as a row-stochastic matrix to mix values:
\begin{equation}
  H \;=\; A V \;\in\; \R^{n \times \dhead}.
\end{equation}
Row \(t\) of \(H\) is the weighted sum \(\sum_s A_{t,s} v_s\) of all values,
weighted by how much token \(t\) wanted to attend to token \(s\).

\keyidea{Self-attention is one matrix multiply (\(QK\T\)), one row-wise softmax,
one more matrix multiply (\(AV\)). The three steps look small. The middle
matrix \(A\) is \(n \times n\). That is the entire story of this book.}

\subsection{Multi-head attention}

The single-head version has limited capacity: every position attends to every
other through a single dot-product space of dimension \(\dhead\). In
practice, you run \(\heads\) of these in parallel:
\begin{equation}
  H_h \;=\; \softmax\!\left( \frac{Q_h K_h\T}{\sqrt{\dhead}} \right) V_h,
  \quad h = 1, \ldots, \heads,
\end{equation}
with \(\dhead = d / \heads\). Concatenate the head outputs along the
feature axis and apply a final output projection \(W_O \in \R^{d \times d}\):
\begin{equation}
  \text{MHA}(X) \;=\; [H_1, H_2, \ldots, H_\heads]\, W_O \;\in\; \R^{n \times d}.
\end{equation}
Different heads can specialise; one might track syntax, another co-reference,
another long-range topical similarity. The cost adds linearly in \(\heads\)
because each head shrinks \(\dhead\) proportionally.\sidenote{Concretely, for
\(d = 4096\) and \(\heads = 32\), each head has \(\dhead = 128\). This is
a common choice; \(\dhead\) between 64 and 128 has been the sweet spot since
OpenAI's GPT-3 (third-generation Generative Pre-trained Transformer).}

\subsection{The causal mask}

For an \emph{autoregressive} language model, token \(t\) is not allowed to
look at tokens \(> t\). That would be cheating during training (the model
could just copy the next token from the input). The fix is the
\term{causal mask}: before the softmax, set every entry above the diagonal
of \(QK\T\) to \(-\infty\). The softmax of \(-\infty\) is zero, so
\(A_{t,s} = 0\) whenever \(s > t\). The lower triangular structure is
what the entire generation loop depends on.

\subsection{Why self-attention is genuinely different from an LSTM}

In an LSTM, information from token \(s\) reaches token \(t\) (with \(s \ll t\))
through a chain of \(t - s\) updates of the hidden state. Each update can
distort, drop, or blur the signal.

In self-attention, the contribution of token \(s\) to token \(t\) is a single
dot product followed by a softmax-weighted sum. There is no chain. Distance
in the sequence does not, by itself, attenuate the signal.

\keyidea{A transformer can ``copy'' a fact from token 1 into token 1000 in
a single attention layer. An LSTM has to relay it through 999 hidden-state
updates. This is why transformers are so much better at long-range
dependencies, and why the resulting compute bill is paid in an \(n \times n\)
matrix.}

\subsection{The shapes, summarised}

For a transformer with \(n\) tokens, \(\dmodel = d\), \(\heads\) heads, and
\(\dhead = d / \heads\):

\begin{center}
\begin{tabular}{lll}
\toprule
Tensor & Shape & Notes \\
\midrule
\(X\)   & \([n, d]\)              & input to the block \\
\(Q, K, V\) per head & \([n, \dhead]\) & queries, keys, values \\
\(QK\T\) per head & \([n, n]\)    & \emph{the} matrix \\
\(A\) per head    & \([n, n]\)    & after row-wise softmax \\
\(AV\) per head   & \([n, \dhead]\) & weighted values \\
concat            & \([n, d]\)    & all heads stitched \\
output            & \([n, d]\)    & after \(W_O\) \\
\bottomrule
\end{tabular}
\end{center}

The bold row to remember is \([n, n]\). One \(n \times n\) tensor per head per
layer. For \(n = 10^6\) and 32 heads and 32 layers, that is
\(32 \cdot 32 \cdot (10^6)^2 = 10^{15}\) cells per forward pass. Even at one
byte per cell that is a petabyte of intermediate state. The next chapter
spells out the consequences.

\iffalse
\begin{lstlisting}[caption={Single-head scaled dot-product attention (equation~\ref{eq:attn-mat}), causal LM variant.}]
import math
import torch
import torch.nn.functional as F

def sdpa_single_head(x, W_q, W_k, W_v, causal=True):
    # x: [n, d]; W_q,W_k,W_v: [d, d_head] -> output [n, d_head]
    n, d = x.shape
    d_head = W_q.shape[1]
    q = x @ W_q                     # [n, d_head]
    k = x @ W_k                     # [n, d_head]
    v = x @ W_v                     # [n, d_head]
    scores = (q @ k.T) / math.sqrt(d_head)   # [n, n]
    if causal:
        mask = torch.triu(torch.ones(n, n, dtype=torch.bool), diagonal=1)
        assert mask.shape == (n, n)
        scores = scores.masked_fill(mask, float("-inf"))
    attn = F.softmax(scores, dim=-1)         # row-wise; rows sum to 1
    return attn @ v                           # [n, d_head]
\end{lstlisting}
\fi

\noindent
This is exactly ``\(QK^{\mathsf T}\), mask, softmax, multiply by \(V\)'' with
no fused kernels; production training stacks heads and batches \(B\) leading
axes.

\paragraph{What does an attention head actually learn?}
Interpretability work on small transformers shows heads are not interchangeable:
many behave like fixed positional routes or shallow circuits rather than
generic similarity.\sidenote{See \emph{e.g.}\ ``What Does BERT Look At?''
\citep{clark2019whatbert} for encoder self-attention (BERT: Bidirectional Encoder Representations from Transformers) and Olsson et al.\ on
\emph{induction heads} in small language models (LMs).} A canonical decoder example is an
\emph{induction head}: when the context repeats a pattern \texttt{A B \ldots\
A}, this head attends from the second \texttt{A} back to \texttt{B} so the
model can predict the token that followed the first \texttt{A} --- supporting
in-context copying and small algorithmic steps. The attention matrix is
sparse and structured (previous-token + source match), not a diffuse softmax
over everything. Multi-head attention exists precisely so several such
low-rank circuits can coexist: some heads gate syntax-like agreement, others
implement retrieval-like jumps, others approximate positional offsets. Width
\(d\) split into \(\heads\) heads trades parameter efficiency for a \emph{bank}
of parallel routing patterns.
